\section{Landasan Teori}

Landasan teori pada bagian ini membahas konsep yang digunakan untuk memahami rancangan sistem \textit{checkout} otomatis Nasi Padang. Pembahasan diawali dengan Nasi Padang untuk menjelaskan karakteristik makanan dan bentuk penyajian yang menjadi ruang lingkup penelitian. Selanjutnya, pembahasan mencakup \textit{Computer Vision} dan citra digital, \textit{instance segmentation}, YOLO11-seg, penghitungan instans dan perhitungan total harga, serta evaluasi kinerja model dan sistem. Detail mengenai dataset, konfigurasi pelatihan, perangkat, dan prosedur pengujian dibahas pada Bab III.

\subsection{\textit{Computer Vision} dan Citra Digital}

\textit{Computer Vision} digunakan untuk memperoleh informasi dari citra melalui proses komputasi. Pada analisis makanan, informasi visual dapat dimanfaatkan untuk melakukan pengenalan, deteksi, dan segmentasi objek sehingga jenis serta lokasi makanan pada citra dapat ditentukan \cite{aguilar2017}.

Citra digital berwarna tersusun atas nilai piksel pada beberapa kanal warna. Informasi tersebut menjadi masukan bagi model untuk membentuk representasi fitur yang berkaitan dengan karakteristik visual objek, seperti bentuk, warna, tekstur, dan susunan objek pada citra. Pada sistem yang dirancang, citra piring digunakan sebagai masukan bagi YOLO11-seg untuk mengenali lauk pada tingkat instans.

Karena pengambilan citra pada sistem dilakukan dari posisi tetap di atas piring, area objek dapat diamati dengan perspektif yang relatif konsisten. Pendekatan pengambilan citra dari arah atas juga digunakan pada sistem \textit{visual checkout} makanan untuk memperoleh tampilan objek pada area pengamatan secara konsisten \cite{wu2021}.

\subsection{\textit{Instance Segmentation}}

\textit{Instance segmentation} merupakan tugas \textit{Computer Vision} yang menghasilkan informasi kelas dan wilayah piksel untuk setiap instans objek. Berbeda dengan segmentasi semantik yang memberikan label kelas pada piksel tanpa membedakan objek dari kelas yang sama, \textit{instance segmentation} mempertahankan identitas setiap objek sebagai instans yang terpisah \cite{he2017maskrcnn}.

Pemisahan pada tingkat instans sesuai dengan kebutuhan sistem karena satu piring dapat mengandung lebih dari satu lauk dari kelas yang sama. Setiap objek perlu dipertahankan sebagai instans yang berbeda agar jumlahnya dapat dihitung secara individual. Pemanfaatan segmentasi dan penghitungan pada tingkat instans juga digunakan pada penelitian SibNet dan FoodMask untuk mengenali serta menghitung objek makanan \cite{nguyen2022sibnet,nguyen2023foodmask}.

Kondisi oklusi dapat menjadi tantangan dalam \textit{instance segmentation} karena sebagian area objek dapat tertutup oleh objek lain sehingga informasi visual yang tersedia berkurang \cite{ke2021}. Oleh karena itu, kondisi oklusi dipandang sebagai salah satu faktor yang dapat memengaruhi hasil prediksi model dan bukan sebagai kondisi yang diasumsikan selalu dapat ditangani secara sempurna.

Pada sistem yang dirancang, keluaran \textit{instance segmentation} digunakan untuk memperoleh kelas dan identitas setiap instans lauk. Luas \textit{segmentation mask} tidak digunakan untuk menentukan harga maupun jumlah lauk; jumlah ditentukan berdasarkan banyaknya instans yang dipertahankan sebagai hasil prediksi akhir.

\subsection{YOLO11-seg}

YOLO11-seg merupakan varian YOLO11 yang mendukung tugas \textit{instance segmentation}. Model ini digunakan untuk menghasilkan prediksi lokasi objek, kelas, dan \textit{segmentation mask} dalam satu model sehingga proses klasifikasi lauk dan segmentasi tidak memerlukan dua model yang terpisah \cite{khanam2024,qiu2025}.

\subsubsection*{a. Arsitektur YOLO11-seg}

Secara umum, arsitektur YOLO11 terdiri atas tiga bagian utama, yaitu \textit{backbone}, \textit{neck}, dan \textit{head}. \textit{Backbone} mengekstraksi fitur dari citra masukan, \textit{neck} menggabungkan fitur dari beberapa tingkat resolusi, sedangkan \textit{head} menggunakan fitur tersebut untuk menghasilkan prediksi objek \cite{khanam2024}. Pada YOLO11-seg, bagian keluaran juga mencakup informasi yang diperlukan untuk membentuk segmentasi setiap instans \cite{qiu2025}.

Pada bagian \textit{backbone}, citra diproses melalui lapisan konvolusi dan modul utama seperti C3k2, SPPF (\textit{Spatial Pyramid Pooling--Fast}), dan C2PSA. C3k2 digunakan dalam proses ekstraksi fitur, SPPF memperluas informasi konteks melalui operasi \textit{pooling}, sedangkan C2PSA memberikan mekanisme perhatian terhadap fitur yang diproses \cite{khanam2024,qiu2025}.

Fitur dari \textit{backbone} diteruskan ke bagian \textit{neck}. Bagian ini menggabungkan informasi dari beberapa tingkat resolusi melalui operasi seperti \textit{upsampling} dan \textit{concatenation}. Penggabungan fitur lintas skala memungkinkan informasi spasial dari fitur beresolusi lebih tinggi dipadukan dengan informasi semantik dari fitur yang lebih dalam \cite{khanam2024,qiu2025}.

Hasil dari \textit{neck} kemudian diteruskan ke \textit{head}. Pada YOLO11-seg, bagian deteksi menghasilkan prediksi kelas dan posisi \textit{bounding box}, sedangkan bagian segmentasi menghasilkan informasi \textit{mask} untuk setiap kandidat objek. Salah satu mekanisme yang digunakan adalah pembentukan \textit{prototype mask} yang dikombinasikan dengan koefisien \textit{mask} setiap objek untuk menghasilkan \textit{mask} masing-masing instans \cite{qiu2025}. Penerapan YOLO11-seg pada penelitian lain juga menunjukkan penggunaan C3k2, SPPF, C2PSA, penggabungan fitur lintas skala, serta keluaran segmentasi dalam satu arsitektur \cite{syafira2025}.

\subsubsection*{b. Pelatihan YOLO11-seg}

Pelatihan model dapat dimulai dari bobot pralatih agar fitur visual yang telah dipelajari sebelumnya dapat digunakan sebagai titik awal dan kemudian disesuaikan dengan kelas pada dataset penelitian. Pada penelitian ini, bobot YOLO11-seg disesuaikan untuk mengenali sembilan kelas, yaitu nasi, rendang, ayam goreng, perkedel, sayur nangka, daun singkong, telur dadar, telur kari, dan sambal.

Pada proses pelatihan, citra dan \textit{ground truth} diberikan kepada model untuk menghasilkan prediksi lokasi objek, kelas, dan segmentasi. Prediksi tersebut dibandingkan dengan \textit{ground truth} untuk memperoleh nilai \textit{loss}. Nilai \textit{loss} kemudian digunakan dalam proses \textit{backpropagation} untuk memperbarui bobot model secara berulang. Data validasi digunakan untuk memantau performa model selama pelatihan, sedangkan data pengujian digunakan setelah model akhir diperoleh.

Nilai parameter pelatihan, seperti varian YOLO11-seg, ukuran citra masukan, jumlah \textit{epoch}, \textit{batch size}, konfigurasi augmentasi, dan parameter lain yang digunakan pada implementasi dijelaskan pada Bab III.

\subsection{Penghitungan Instans}

Penghitungan lauk dilakukan berdasarkan hasil prediksi akhir YOLO11-seg. Setiap prediksi akhir merepresentasikan satu instans objek, sedangkan label kelas menentukan kelompok tempat instans tersebut dihitung. Dengan demikian, penghitungan dilakukan dari jumlah instans yang dikenali dan bukan dari jumlah piksel atau luas \textit{segmentation mask}.

Program memeriksa hasil prediksi dan menambahkan jumlah pada kelas yang sesuai untuk setiap instans yang ditemukan. Apabila hasil inferensi menghasilkan dua instans rendang dan satu instans perkedel, maka jumlah rendang adalah dua dan jumlah perkedel adalah satu. Prinsip penghitungan pada tingkat instans juga digunakan pada penelitian penghitungan makanan seperti SibNet dan FoodMask \cite{nguyen2022sibnet,nguyen2023foodmask}.

Informasi jumlah setiap kelas selanjutnya menjadi masukan bagi proses kalkulasi harga. Pemisahan antara proses pengenalan visual dan proses penghitungan memungkinkan model tetap bertugas mengenali objek, sedangkan program melakukan agregasi jumlah berdasarkan label kelas hasil prediksi.

\subsection{Penghitungan Instans dan Perhitungan Total Harga}

Penghitungan objek makanan dilakukan berdasarkan hasil prediksi akhir YOLO11-seg. Setiap prediksi akhir merepresentasikan satu instans objek, sedangkan label kelas menentukan kelompok tempat instans tersebut dihitung. Dengan demikian, penghitungan dilakukan berdasarkan jumlah instans yang dikenali, bukan berdasarkan jumlah piksel atau luas \textit{segmentation mask}.

Program memeriksa setiap hasil prediksi dan menambahkan jumlah pada kelas yang sesuai. Apabila hasil inferensi menghasilkan dua instans rendang dan satu instans perkedel, jumlah rendang yang tercatat adalah dua dan jumlah perkedel adalah satu. Prinsip penghitungan pada tingkat instans juga digunakan dalam penelitian penghitungan makanan seperti SibNet dan FoodMask \cite{nguyen2022sibnet,nguyen2023foodmask}.

Jumlah instans pada setiap kelas selanjutnya digunakan untuk menghitung total harga. Setiap kelas memiliki harga satuan tetap yang disimpan sebagai data pada program. YOLO11-seg tidak menentukan harga makanan; model hanya menghasilkan kelas dan instans objek, sedangkan proses penghitungan harga dilakukan oleh program. Pemanfaatan hasil pengenalan citra sebagai dasar perhitungan harga juga diterapkan pada penelitian mengenai \textit{automatic billing} dan \textit{visual checkout} makanan \cite{wu2021,guo2023}.

Subtotal suatu kelas diperoleh dengan mengalikan jumlah instans pada kelas tersebut dengan harga satuannya. Total harga kemudian diperoleh dengan menjumlahkan subtotal dari seluruh kelas yang dikenali. Perhitungan tersebut dinyatakan pada Persamaan~\ref{eq:total-harga}.

\begin{equation}
    T = \sum_{i=1}^{C} n_i p_i
    \label{eq:total-harga}
\end{equation}

dengan:
\begin{itemize}
    \item $T$ adalah total harga yang harus dibayarkan;
    \item $C$ adalah jumlah kelas makanan;
    \item $n_i$ adalah jumlah instans yang dikenali pada kelas ke-$i$; dan
    \item $p_i$ adalah harga satuan makanan pada kelas ke-$i$.
\end{itemize}

Sebagai contoh, apabila sistem mengenali dua instans rendang dan satu instans perkedel, subtotal rendang diperoleh dengan mengalikan dua dengan harga satuan rendang. Subtotal perkedel diperoleh dengan mengalikan satu dengan harga satuan perkedel. Total pembayaran merupakan penjumlahan kedua subtotal tersebut beserta subtotal kelas lain yang dikenali.

Penentuan harga tidak menggunakan berat, volume, luas \textit{mask}, maupun ukuran porsi. Harga satuan disimpan secara terpisah dari model sehingga perubahan harga tidak memerlukan pelatihan ulang YOLO11-seg selama identitas kelas pada program tetap sesuai dengan keluaran model.

\subsection{Evaluasi Kinerja Model dan Sistem}

Evaluasi dilakukan untuk menilai kemampuan YOLO11-seg dalam mengenali objek serta kemampuan sistem dalam mengolah hasil prediksi menjadi jumlah lauk dan total harga. Evaluasi model mencakup hasil deteksi dan segmentasi, sedangkan evaluasi sistem mencakup ketepatan penghitungan instans dan kecepatan pemrosesan.

\subsubsection*{a. Evaluasi Deteksi dan Segmentasi}

Kesesuaian antara prediksi dan \textit{ground truth} dapat diukur menggunakan \textit{Intersection over Union} (IoU). IoU merupakan perbandingan antara luas irisan dan luas gabungan wilayah prediksi dengan \textit{ground truth}, dan digunakan dalam evaluasi model segmentasi berbasis YOLO \cite{qiu2025}. Jika $A$ merupakan wilayah prediksi dan $B$ merupakan wilayah \textit{ground truth}, IoU dinyatakan sebagai:

\begin{equation}
    IoU(A,B) = \frac{|A \cap B|}{|A \cup B|}
    \label{eq:iou}
\end{equation}

Berdasarkan pencocokan antara prediksi dan \textit{ground truth}, hasil prediksi dapat dikelompokkan menjadi \textit{true positive} (TP), \textit{false positive} (FP), dan \textit{false negative} (FN). \textit{Precision}, \textit{recall}, dan \textit{F1-score} merupakan metrik yang umum digunakan untuk mengevaluasi hasil deteksi dan \textit{instance segmentation}. Rumus ketiga metrik tersebut juga digunakan dalam penelitian \textit{instance segmentation} berbasis YOLOv8-Seg \cite{yue2023}.

\textit{Precision} menunjukkan proporsi prediksi positif yang benar dan dinyatakan sebagai:

\begin{equation}
    \textit{Precision} = \frac{TP}{TP + FP}\times 100\%
    \label{eq:precision}
\end{equation}

\textit{Recall} menunjukkan proporsi objek sebenarnya yang berhasil dikenali dan dinyatakan sebagai:

\begin{equation}
    \textit{Recall} = \frac{TP}{TP + FN}\times 100\%
    \label{eq:recall}
\end{equation}

Nilai \textit{F1-score} merupakan rata-rata harmonik dari \textit{precision} dan \textit{recall}, yang dinyatakan sebagai:

\begin{equation}
    F1 = 2 \times
		{\textit{Precision}} \times
    \frac{\textit{Recall}}
         {\textit{Precision} + \textit{Recall}}
    \label{eq:f1}
\end{equation}

Selain itu, performa YOLO11-seg dinilai menggunakan \textit{mean Average Precision} (mAP). mAP50 menunjukkan nilai mAP pada ambang IoU 0,50, sedangkan mAP50--95 merangkum performa pada rentang ambang IoU 0,50 hingga 0,95. Pada model segmentasi, metrik mAP dapat dilaporkan untuk hasil \textit{bounding box} dan \textit{segmentation mask} \cite{qiu2025,syafira2025}.

\subsubsection*{b. Evaluasi Penghitungan Instans}

Metrik deteksi dan segmentasi belum secara langsung menunjukkan seberapa tepat jumlah objek yang dihasilkan sistem. Pada penelitian SibNet, \textit{Mean Absolute Error} (MAE) digunakan untuk mengevaluasi kesalahan penghitungan instans makanan berdasarkan selisih absolut antara jumlah prediksi dan jumlah acuan \cite{nguyen2022sibnet}. Pada penelitian ini, bentuk MAE tersebut disesuaikan untuk menghitung kesalahan penghitungan pada masing-masing kelas lauk. Untuk kelas ke-$i$ pada $N$ citra pengujian, MAE dinyatakan dengan Persamaan~\ref{eq:mae-count}.

\begin{equation}
    MAE_i = \frac{1}{N} \sum_{j=1}^{N} \left|\hat{n}_{ji} - n_{ji}\right|
    \label{eq:mae-count}
\end{equation}

Pada persamaan tersebut, $i$ menunjukkan kelas yang sedang dinilai, $j$ menunjukkan citra pengujian, $\hat{n}_{ji}$ merupakan jumlah instans hasil prediksi pada citra ke-$j$ untuk kelas ke-$i$, sedangkan $n_{ji}$ merupakan jumlah instans berdasarkan \textit{ground truth}. Nilai MAE yang lebih kecil menunjukkan rata-rata kesalahan penghitungan yang lebih rendah, sedangkan nilai nol menunjukkan tidak terdapat selisih jumlah pada data yang dievaluasi.

\subsubsection*{c. Kecepatan Inferensi}

Kecepatan model dapat dinilai berdasarkan waktu inferensi rata-rata yang diperlukan untuk memproses satu citra. Penggunaan waktu inferensi sebagai salah satu indikator performa juga ditemukan pada penelitian segmentasi berbasis YOLO untuk menilai kesesuaian model terhadap kebutuhan pemrosesan secara \textit{real-time} \cite{syafira2025}. Perangkat, data pengujian, dan prosedur pengukuran yang digunakan dijelaskan pada Bab III.

